\chapter{初识分数}\label{ch:g4:fractions}

半个苹果，四分之一小时，全班的四分之三：
分数用来表示事物的一部分。本章介绍
$\frac{a}{b}$ 的写法，学习在图形和数轴上读出分数，
并比较最简单的分数。（分数将在
\cref{ch:g6:fractions}中继续学习。）

\section{为各部分命名}

\begin{definition}[分数]\label{def:g4:fractions:def}
把一个整体平均分成若干等份，取其中几份 --- 这就是一个
\emph{分数}\index{分数}：
\[
\frac{3}{4}
\quad
\begin{array}{l}
3 \text{：取了几份（\emph{分子}\index{分子}）；}\\
4 \text{：整体被平均分成几份（
\emph{分母}\index{分母}）。}
\end{array}
\]
$\frac12$ 是\emph{一半}，$\frac13$ 是\emph{三分之一}，$\frac14$ 是
\emph{四分之一}，$\frac{1}{10}$ 是\emph{十分之一}。
\end{definition}

\begin{omfigure}
\begin{tikzpicture}[scale=0.78]
  % half circle
  \draw[thick] (0,0) circle (1.0);
  \fill[omDef!40] (0,0) -- (1,0) arc (0:180:1) -- cycle;
  \draw[thick] (-1,0) -- (1,0);
  \node[below, font=\small] at (0,-1.2) {$\frac12$};
  % quarters
  \begin{scope}[xshift=3.6cm]
  \draw[thick] (0,0) circle (1.0);
  \fill[omThm!40] (0,0) -- (1,0) arc (0:90:1) -- cycle;
  \fill[omThm!40] (0,0) -- (0,1) arc (90:180:1) -- cycle;
  \fill[omThm!40] (0,0) -- (-1,0) arc (180:270:1) -- cycle;
  \draw[thick] (-1,0) -- (1,0);
  \draw[thick] (0,-1) -- (0,1);
  \node[below, font=\small] at (0,-1.2) {$\frac34$};
  \end{scope}
  % thirds bar
  \begin{scope}[xshift=6.6cm, yshift=-0.35cm]
  \foreach \i in {0,1,2} \draw[omExo] (\i*1.3,0) rectangle
    (\i*1.3+1.3,0.7);
  \fill[omMeth!40] (0.04,0.04) rectangle (1.26,0.66);
  \draw[thick] (0,0) rectangle (3.9,0.7);
  \node[below, font=\small] at (1.95,-0.15) {$\frac13$};
  \end{scope}
  % tenths bar
  \begin{scope}[xshift=11.6cm, yshift=-0.35cm]
  \foreach \i in {0,...,9} \draw[omExo] (\i*0.42,0) rectangle
    (\i*0.42+0.42,0.7);
  \foreach \i in {0,...,6} \fill[omProp!40] (\i*0.42+0.03,0.04)
    rectangle (\i*0.42+0.39,0.66);
  \draw[thick] (0,0) rectangle (4.2,0.7);
  \node[below, font=\small] at (2.1,-0.15) {$\frac{7}{10}$};
  \end{scope}
\end{tikzpicture}

{\small 在图形上读分数。各部分必须
\emph{相等}：三块大小不等的部分不能算作三分之一！}
\end{omfigure}

\begin{example}[一组物体的分数]\label{ex:g4:fractions:collection}
$\frac13$ 的 $12$ 颗弹珠：把这 $12$ 颗弹珠平均分成 $3$
组，每组 $4$ 颗；一组是 $\frac13$，所以 $\frac13$ 的 $12$ 是 $4$
--- 而 $\frac23$ 的 $12$ 颗是两组，即 $8$ 颗。
\end{example}

\section{在数轴上认识分数}

\begin{method}[在数轴上表示 $\frac{a}{b}$]\label{met:g4:fractions:line}
把数轴上每个单位长度平均分成 $b$ 份，从 $0$ 出发，走 $a$
小段。若 $a$ 大于 $b$，则走过的路线越过 $1$：分数可以
比一个整体还大！
\end{method}

\begin{omfigure}
\begin{tikzpicture}[scale=1.6]
  \draw[-Stealth] (-0.2,0) -- (6.9,0);
  \foreach \i in {0,4,8,12,16,20,24}
    \draw ({\i*0.25},-0.09) -- ({\i*0.25},0.09);
  \foreach \i in {0,...,24} \draw ({\i*0.25},-0.05) -- ({\i*0.25},0.05);
  \foreach \x/\l in {0/0, 1/1, 2/2, 3/3, 4/4, 5/5, 6/6}
    \node[below, font=\small] at (\x,-0.1) {$\l$};
  \fill[omDef] (0.5,0) circle (1.6pt)
    node[above, font=\small] {$\frac24$};
  \fill[omThm] (1.25,0) circle (1.6pt)
    node[above, font=\small] {$\frac54$};
  \fill[omMeth] (2.75,0) circle (1.6pt)
    node[above, font=\small] {$\frac{11}{4}$};
\end{tikzpicture}

{\small 数轴被四等分：$\frac24$ 与
$\frac12$ 落在同一点；$\frac54$ 在 $1$ 右侧一个四分之一处；$\frac{11}{4}$ 是
$2$ 与 $\frac34$。}
\end{omfigure}

\begin{example}[藏在分数里的整数]\label{ex:g4:fractions:whole}
$\frac44 = 1$（四个四分之一合成一个整体）；$\frac82 = 4$；
$\frac{12}{3} = 4$。而 $\frac74$ 是 $1 + \frac34$：一个整体
再加三个四分之一。
\end{example}

\section{比较与相加简单分数}

\begin{proposition}[分母相同]\label{prop:g4:fractions:compare}
分母相同时，各份的大小相同 --- 因此只要
数出份数即可：
\[
\frac{2}{8} < \frac{5}{8},
\qquad
\frac{3}{8} + \frac{4}{8} = \frac{7}{8}.
\]
与一个整体比较也很容易：$\frac{5}{8} < 1 < \frac{9}{8}$
（一个少于、一个多于 $8$ 个八分之一）。
\end{proposition}

\begin{proof}[借助图形说明理由]
在同一根长条上先涂 $3$ 个八分之一，再涂 $4$ 个八分之一，
共涂出 $7$ 个八分之一。
\end{proof}

\begin{example}[分母不同，点相同]\label{ex:g4:fractions:equal}
在数轴上，$\frac12$、$\frac24$ 和 $\frac{5}{10}$ 都落在
同一个点上：它们是同一个数的不同写法。把每一半
平均分成两份，得到四分之一；平均分成五份，得到十分之一。（一般
规律在\cref{ch:g6:fractions}中给出。）
\end{example}

\begin{example}[一小时的四分之一]\label{ex:g4:fractions:hour}
一小时有 $60$ 分钟，所以一小时的四分之一是
$60 \div 4 = 15$ 分钟，一小时的四分之三是
$45$ 分钟。“两点半”是
两点过 $\frac12$ 小时。
\end{example}

\section{练习}

\begin{exercise}[$\star$]\label{exo:g4:fractions:1}
写出所示的分数：一个披萨平均切成 $6$ 块，还剩 $5$ 块；
一块有 $8$ 小格的巧克力，吃掉了 $3$ 小格（写出吃掉的分数）；一面
旗子竖着平均分成 $3$ 等份，涂色的有 $1$ 份。
\end{exercise}

\begin{exercise}[$\star$]\label{exo:g4:fractions:2}
画一根长条，平均分成 $5$ 份，涂出 $\frac35$。再在另一根同样长的长条上
涂出 $\frac{2}{5}$。哪个涂色部分更大？
\end{exercise}

\begin{exercise}[$\star$]\label{exo:g4:fractions:3}
计算：$\frac12$ 乘 $18$；\quad $\frac14$ 乘 $20$；\quad
$\frac13$ 乘 $21$；\quad $\frac34$ 乘 $20$（三组
四分之一）。
\end{exercise}

\begin{exercise}[$\star$]\label{exo:g4:fractions:4}
在三等分的数轴上标出：$\frac13$；\ $\frac33$；\
$\frac53$；\ $2$。哪个分数恰好落在 $2$ 上？
\end{exercise}

\begin{exercise}[$\star$]\label{exo:g4:fractions:5}
抄写并在空处填上 $<$、$>$ 或 $=$：
\[
\frac38 \;?\; \frac68, \qquad
\frac55 \;?\; 1, \qquad
\frac74 \;?\; 1, \qquad
\frac12 \;?\; \frac24 .
\]
\end{exercise}

\begin{exercise}[$\star$]\label{exo:g4:fractions:6}
计算：$\frac26 + \frac36$；\quad $\frac58 - \frac28$；\quad
$\frac14 + \frac24$；\quad $1 - \frac13$（几个三分之一
等于 $1$？）
\end{exercise}

\begin{exercise}[$\star$]\label{exo:g4:fractions:7}
写成一个整数加一个小于 $1$ 的分数的形式：
$\frac54$; \quad $\frac73$; \quad $\frac{9}{2}$.
\end{exercise}

\begin{exercise}[$\star$]\label{exo:g4:fractions:8}
下列时间各是多少分钟：一小时的一半；一小时的四分之一；
一小时的四分之三；一小时的 $\frac13$？
\end{exercise}

\begin{exercise}[$\star$]\label{exo:g4:fractions:9}
一个班有 $24$ 名学生，其中 $\frac14$ 的学生戴眼镜。戴眼镜的
有多少名？不戴眼镜的有多少名？
\end{exercise}

\begin{exercise}[$\star\star$]\label{exo:g4:fractions:10}
尼娜吃了一个披萨的 $\frac38$，萨姆吃了同一个
披萨的 $\frac48$。他们共吃了几分之几？还剩几分之几？
谁吃得多？
\end{exercise}

\begin{exercise}[$\star\star$]\label{exo:g4:fractions:11}
判断下列说法是否正确，并借助图画或例子说明理由：“小披萨的 $\frac12$ 少于
很大披萨的 $\frac14$ 是可能的”。
比较同一个整体的两个分数之前，
必须始终知道什么？
\end{exercise}
